\documentclass[11pt,a4paper]{article}
% Preámbulo: tipografía, paleta monocromática y macros de notación.
\usepackage[T1]{fontenc}
\usepackage[utf8]{inputenc}
\usepackage{lmodern}
\usepackage{amsmath}
\usepackage[spanish,es-noquoting,es-noshorthands,es-tabla]{babel}
\usepackage[a4paper,margin=2.6cm,headsep=14pt]{geometry}
\usepackage{microtype}
\usepackage{booktabs}
\usepackage{longtable}
\usepackage{array}
\usepackage{xcolor}
\usepackage{fancyhdr}
\usepackage{titlesec}
\usepackage{enumitem}
\usepackage{fancyvrb}
\usepackage{tikz}
\usetikzlibrary{automata,positioning,arrows.meta,calc}
\usepackage[hidelinks,bookmarks=true]{hyperref}

% --- Paleta monocromática: sólo grises, con un único acento. -----------------
\definecolor{tinta}{gray}{0.13}
\definecolor{grisMedio}{gray}{0.42}
\definecolor{grisSuave}{gray}{0.70}
\definecolor{grisFondo}{gray}{0.955}
\definecolor{acento}{HTML}{2F4858}

\color{tinta}

% --- Títulos ----------------------------------------------------------------
\titleformat{\section}{\Large\bfseries\color{acento}}{\thesection}{0.7em}{}
\titleformat{\subsection}{\large\bfseries\color{tinta}}{\thesubsection}{0.7em}{}
\titleformat{\subsubsection}{\bfseries\color{grisMedio}}{\thesubsubsection}{0.7em}{}
\titlespacing*{\section}{0pt}{2.2ex plus 1ex minus .2ex}{1.2ex plus .2ex}

% --- Encabezado y pie -------------------------------------------------------
\setlength{\headheight}{14pt}
\pagestyle{fancy}
\fancyhf{}
\renewcommand{\headrulewidth}{0.4pt}
\renewcommand{\headrule}{\hbox to\headwidth{\color{grisSuave}\leaders\hrule height \headrulewidth\hfill}}
\fancyhead[L]{\footnotesize\color{grisMedio}Compiladores e Intérpretes — Etapa 1}
\fancyhead[R]{\footnotesize\color{grisMedio}\nouppercase{\leftmark}}
\fancyfoot[C]{\footnotesize\color{grisMedio}\thepage}

% --- Bloques de código ------------------------------------------------------
\DefineVerbatimEnvironment{codigo}{Verbatim}%
  {fontsize=\small,xleftmargin=1.1em,commandchars=\\\{\}}
\DefineVerbatimEnvironment{codigoLlano}{Verbatim}%
  {fontsize=\small,xleftmargin=1.1em}

% --- Macros de notación -----------------------------------------------------
\newcommand{\bs}{\textbackslash}
\newcommand{\caret}{\textasciicircum}
\newcommand{\lex}[1]{\texttt{#1}}
\newcommand{\tok}[1]{\textnormal{\textsf{\small #1}}}
% --- Cardinalidad como exponente propiamente dicho -------------------------
% \sups{c}    exponente suelto:  ( ... )^*
% \ere{X}{c}  símbolo con su exponente:  Digito^{1-9}
\newcommand{\sups}[1]{\textsuperscript{\normalfont\texttt{\footnotesize #1}}}
\newcommand{\ere}[2]{\texttt{#1}\sups{#2}}
\newcommand{\rep}[2]{\sups{#1--#2}}
\newcommand{\alt}{\ensuremath{\;|\;}}

% Nota al pie de notación, repetible: se recuerda la convención donde aparece.
\newcommand{\notaKleene}{\footnote{\textbf{Notación.} \ere{X}{*} = cero o más
repeticiones de \lex{X} (estrella de Kleene); \ere{X}{+} = una o más;
\ere{X}{?} = cero o una; \ere{X}{1--9} = entre 1 y 9 repeticiones, \emph{ambos
extremos inclusive}; \lex{A | B} = alternativa. La tabla completa está en la
sección~\ref{sec:notacion}.}}

\newenvironment{er}{%
  \par\smallskip\begingroup\ttfamily\small\parindent=0pt
  \leftskip=1.2em\obeylines}{\par\endgroup\smallskip}

\setlist[itemize]{leftmargin=1.3em,itemsep=0.25ex,topsep=0.5ex}
\setlist[description]{leftmargin=1.3em,itemsep=0.35ex,topsep=0.6ex}
\renewcommand{\arraystretch}{1.25}


\begin{document}

% ===========================================================================
% Portada
% ===========================================================================
\begin{titlepage}
\thispagestyle{empty}
\raggedright
\vspace*{2.2cm}

{\color{grisMedio}\footnotesize\sffamily
UNIVERSIDAD NACIONAL DEL SUR\\
Departamento de Ciencias e Ingeniería de la Computación\\
Compiladores e Intérpretes --- Segundo Cuatrimestre de 2026\par}

\vspace{2.4cm}
{\color{grisSuave}\rule{\textwidth}{0.6pt}}
\vspace{1.1cm}

{\Huge\bfseries\color{acento} Compilador de MiniJava\par}
\vspace{0.5cm}
{\LARGE\color{tinta} Etapa 1 --- Analizador Léxico\\[0.25cm]
\large\color{grisMedio} y Módulo Principal\par}

\vspace{1.1cm}
{\color{grisSuave}\rule{\textwidth}{0.6pt}}

\vspace{2.6cm}
{\large\color{tinta} Agustín Paul\par}
\vspace{0.35cm}
{\color{grisMedio} Informe de entrega\par}

\vfill
{\color{grisMedio}\footnotesize
Implementación por \emph{Método por Estado} (patrón State).\\
Los 56 componentes léxicos, sus expresiones regulares y sus autómatas.\par}
\vspace{0.6cm}
{\color{grisMedio}\footnotesize 29 de agosto de 2026\par}
\end{titlepage}

% ===========================================================================
\thispagestyle{empty}
\tableofcontents
\newpage
\setcounter{page}{1}

% ===========================================================================
\section{Alcance}
% ===========================================================================

Este informe documenta la Etapa 1 del compilador de MiniJava: el \textbf{Analizador Léxico} y el
\textbf{Módulo Principal} que lo consume por consola. Cubre la expresión regular de los 56
componentes léxicos sin omitir ninguno, los autómatas que los reconocen, las decisiones tomadas
donde la especificación admitía más de una lectura, el tratamiento de errores y los ocho logros
perseguidos.

\subsection{Qué hace el programa}

Toma un archivo fuente MiniJava y emite por salida estándar la secuencia de sus componentes léxicos,
uno por línea, en el formato \lex{(NombreToken,Lexema,NroLinea)}. Si no hubo errores, cierra con
\lex{[SinErrores]}; si los hubo, emite un reporte por cada uno, terminado en el código de contrato
\lex{[Error:Lexema|NroLinea]}.

\begin{codigoLlano}
$ java -jar Compilador.jar programa1.java
(litString,"hola",1)
(idMV,v1,2)
(op+,+,2)
(idMV,chau,2)
(pr_if,if,3)
(pr_class,class,3)
(llaveC,},3)
(EOF,$,3)

[SinErrores]
\end{codigoLlano}

El renglón en blanco antes de \lex{[SinErrores]} --- y el equivalente antes de
\lex{[Error:\ldots]} --- son los de los ejemplos de \emph{Pautas Etapa 1}. El archivo fuente puede
tener cualquier extensión: el compilador no inspecciona el nombre, sólo abre la ruta que recibe.

\subsection{Cómo compilarlo y ejecutarlo}

\begin{codigoLlano}
cd compilador
mvn package                                   # genera target/Compilador.jar
java -jar target/Compilador.jar programa1.java

mvn test        # suite propia + los 2 testers de la cátedra, en la misma pasada
mvn verify      # además exige los umbrales de cobertura
\end{codigoLlano}

Requiere Java 25 (\lex{maven.compiler.release=25}).

% ===========================================================================
\section{Notación de expresiones regulares}
\label{sec:notacion}
% ===========================================================================

Toda expresión regular de este informe usa la notación siguiente. La cardinalidad se escribe
siempre como \textbf{exponente} del símbolo o del grupo al que afecta. Los dos operadores que más
conviene distinguir son la estrella de Kleene, que admite \emph{cero o más} repeticiones, y la
repetición acotada \ere{X}{m--n}, que es la que necesita el literal entero.

\begin{center}
\begin{tabular}{@{}ll@{}}
\toprule
\textbf{Forma} & \textbf{Significado} \\
\midrule
\ere{X}{*} & cero o más repeticiones de \lex{X} (estrella de Kleene) \\
\ere{X}{+} & una o más repeticiones de \lex{X} \\
\ere{X}{?} & cero o una repetición de \lex{X} \\
\ere{X}{m--n} & entre \lex{m} y \lex{n} repeticiones, \textbf{ambos extremos inclusive} \\
\lex{A | B} & alternativa: \lex{A} o \lex{B} \\
\lex{( \ldots )}\sups{c} & agrupación; el exponente \lex{c} afecta al grupo entero \\
\bottomrule
\end{tabular}
\end{center}

\medskip
Y las clases de caracteres que se usan como abreviatura:

\begin{codigoLlano}
Digito     = [0-9]
LetraMayus = [A-Z]
LetraMin   = [a-z]
Letra      = LetraMayus | LetraMin
Hex        = [0-9a-fA-F]
\end{codigoLlano}

\begin{quote}\small\color{grisMedio}
Una precisión sobre ``letra'' y ``dígito'': la implementación los reconoce con
\lex{Character.isLetter} y \lex{Character.isDigit}, que abarcan más que el ASCII. Las consecuencias
están en la sección~\ref{sec:divergencias}.
\end{quote}

% ===========================================================================
\section{Los 56 componentes léxicos}
% ===========================================================================

\begin{center}
\begin{tabular}{@{}llp{9.2cm}@{}}
\toprule
\textbf{Familia} & \textbf{Cant.} & \textbf{Tokens} \\
\midrule
Palabras clave & 21 & \tok{pr\_class}, \tok{pr\_extends}, \tok{pr\_interface},
  \tok{pr\_implements}, \tok{pr\_static}, \tok{pr\_boolean}, \tok{pr\_char}, \tok{pr\_int},
  \tok{pr\_void}, \tok{pr\_public}, \tok{pr\_if}, \tok{pr\_else}, \tok{pr\_while},
  \tok{pr\_return}, \tok{pr\_var}, \tok{pr\_this}, \tok{pr\_new}, \tok{pr\_null},
  \tok{pr\_true}, \tok{pr\_false}, \tok{pr\_float} \\
Identificadores & 3 & \tok{idClase}, \tok{idGen}, \tok{idMV} \\
Literales & 4 & \tok{intLiteral}, \tok{charLiteral}, \tok{litString}, \tok{floatLiteral} \\
Puntuación & 10 & \tok{parenA}, \tok{parenC}, \tok{llaveA}, \tok{llaveC}, \tok{corcheteA},
  \tok{corcheteC}, \tok{puntoYComa}, \tok{coma}, \tok{punto}, \tok{dosPuntos} \\
Operadores & 17 & \tok{op\_mayor}, \tok{op\_menor}, \tok{op\_negacion}, \tok{op\_asignacion},
  \tok{op\_igual}, \tok{op\_mayorIgual}, \tok{op\_menorIgual}, \tok{op\_distinto}, \tok{op\_and},
  \tok{op\_or}, \tok{op\_modulo}, \tok{op+}, \tok{op-}, \tok{op\_mult}, \tok{op\_div},
  \tok{op\_incremento}, \tok{op\_decremento} \\
Fin de archivo & 1 & \tok{EOF} \\
\midrule
\textbf{Total} & \textbf{56} & \\
\bottomrule
\end{tabular}
\end{center}

\medskip
Los nombres de la columna derecha son literales: son exactamente los que se imprimen, porque el
formato de salida fijado es \lex{(NombreToken,Lexema,NroLinea)}.

\subsection{Identificadores}

Los tres se distinguen por la caja de su primera letra y, entre los dos que empiezan con mayúscula,
por su longitud.\notaKleene

\begin{center}
\begin{tabular}{@{}lll@{}}
\toprule
\textbf{Token} & \textbf{Expresión regular} & \textbf{Ejemplos} \\
\midrule
\tok{idGen} & \lex{LetraMayus} & \lex{A}, \lex{T} \\
\tok{idClase} & \lex{LetraMayus ( Letra | Digito | \_ )}\sups{+} & \lex{Alpha}, \lex{A1}, \lex{Mi\_Clase} \\
\tok{idMV} & \lex{LetraMin ( Letra | Digito | \_ )}\sups{*} & \lex{v1}, \lex{x}, \lex{mi\_metodo} \\
\bottomrule
\end{tabular}
\end{center}

\medskip
\tok{idGen} y \tok{idClase} comparten el primer carácter y se distinguen por longitud: \textbf{una}
mayúscula sola es \tok{idGen}; una mayúscula seguida de al menos un carácter más es \tok{idClase}.
Ahí está la diferencia entre los dos exponentes de la tabla: el \sups{+} de \tok{idClase} exige ese
carácter adicional, mientras que el \sups{*} de \tok{idMV} admite el identificador de una sola
letra.

Un identificador que empieza con minúscula es siempre \tok{idMV}, \textbf{salvo} que su lexema
completo coincida con una de las 21 palabras clave; ver la sección~\ref{sec:tabla}.

\subsection{\tok{intLiteral}}

\begin{er}
intLiteral = \ere{Digito}{1--9}
\end{er}

Entre uno y nueve dígitos, ambos extremos inclusive. La cota superior es parte de la
especificación del lenguaje, no una limitación de la implementación, y por eso la repetición es
\textbf{acotada} y no una estrella de Kleene: \ere{Digito}{*} describiría un lenguaje distinto ---y
además admitiría el literal vacío---. Escrita sin la abreviatura de cardinalidad:

\begin{er}
intLiteral = Digito \ere{Digito}{?} \ere{Digito}{?} \ere{Digito}{?} \ere{Digito}{?} \ere{Digito}{?} \ere{Digito}{?} \ere{Digito}{?} \ere{Digito}{?}
\end{er}

\noindent
El autómata que le corresponde es, por lo tanto, una \textbf{cadena de nueve estados terminales sin
ciclo}: un bucle sobre un único estado aceptaría corridas de cualquier longitud, que es justamente lo
que la cota prohíbe. Cada uno de los nueve estados es de aceptación, porque un literal entero puede
terminar en cualquier cantidad de dígitos entre uno y nueve; el décimo dígito no tiene transición, y
ahí está el error.

\begin{center}
\begin{tikzpicture}[shorten >=1pt,on grid,auto,
    every state/.style={draw=grisMedio,thick,fill=grisFondo,minimum size=0.82cm,font=\scriptsize,
                        inner sep=1pt},
    every edge/.style={draw=grisMedio,->,>={Stealth[round]}},
    every node/.style={font=\scriptsize}]
  \node[state,initial,initial text={}] (i)                     {$e_0$};
  \node[state,accepting]               (d1) [right=2.05cm of i]  {$d_1$};
  \node[state,accepting]               (d2) [right=2.05cm of d1] {$d_2$};
  \node[state,accepting]               (d3) [right=2.05cm of d2] {$d_3$};
  \node                                (p)  [right=1.5cm of d3]  {$\cdots$};
  \node[state,accepting]               (d8) [right=1.5cm of p]   {$d_8$};
  \node[state,accepting]               (d9) [right=2.05cm of d8] {$d_9$};
  \path[->] (i)  edge node {Digito} (d1)
            (d1) edge node {Digito} (d2)
            (d2) edge node {Digito} (d3)
            (d3) edge (p)
            (p)  edge (d8)
            (d8) edge node {Digito} (d9);
\end{tikzpicture}
\end{center}

\paragraph{El límite se valida al final de la corrida, no durante.} El estado consume la corrida
\emph{completa} de dígitos y recién al terminar comprueba que no supere nueve. La razón es que si el
carácter siguiente a la corrida es un \lex{.} o una \lex{e}/\lex{E}, la corrida continúa hacia un
\tok{floatLiteral}, que \textbf{no} tiene ese límite; cortar en el noveno dígito rompería esa
continuación. Por eso \lex{1234567890.5} es un float válido con diez dígitos en la parte entera,
mientras que \lex{1234567890} solo es un error.

\subsection{\tok{charLiteral}}
\label{sec:char}

\begin{er}
charLiteral = ' X '\quad|\quad' \bs\ Y '\quad|\quad' \bs\ U '
\medskip
X = cualquier carácter salvo \bs\ y '
Y = cualquier carácter salvo u
U = escape Unicode
\end{er}

La especificación da aquí una \textbf{producción con dos alternativas}, y las variables \lex{X} e
\lex{Y} son independientes entre sí. Eso hace que la barra invertida y el carácter que la sigue
formen una \emph{unidad}: de ahí que \lex{'\bs\bs'} y \lex{'\bs''} sean literales válidos. El literal
vacío \lex{''} no está en ninguna alternativa y es un error.

Como \lex{Y} es ``cualquier carácter'', \lex{'\bs q'} es válido, con valor \lex{q}: una permisividad
explícita de la especificación, más laxa que Java. Y ``cualquier carácter'' se toma al pie de la
letra: valen las letras acentuadas, la eñe y los alfabetos no latinos igual que las de ASCII. Las dos
únicas excepciones no son restricciones del juego de caracteres sino estructurales ---un salto de
línea o un fin de archivo dejan el literal sin cerrar---, y la \lex{u} es el único carácter que no se
consume como común, porque abre el escape Unicode.

\paragraph{Qué significa exactamente ``un carácter''.} El límite no es ASCII: es que \lex{X} e
\lex{Y} son \textbf{una unidad de código UTF-16}. Viene de la interfaz provista, que declara
\lex{char getNextChar()}, y un \lex{char} de Java son 16 bits.

\begin{center}
\begin{tabular}{@{}llcl@{}}
\toprule
\textbf{Carácter} & \textbf{Code points} & \textbf{Unidades} & \textbf{Resultado} \\
\midrule
eñe & \lex{U+00F1} & 1 & \tok{charLiteral} válido \\
estrella blanca & \lex{U+2B50} & 1 & \tok{charLiteral} válido \\
ideograma CJK & \lex{U+4E2D} & 1 & \tok{charLiteral} válido \\
estrella brillante & \lex{U+1F31F} & 2 (par subrogado) & error: \emph{mal formado} \\
estrella $+$ selector & \lex{U+2B50 U+FE0F} & 2 & error: \emph{mal formado} \\
\bottomrule
\end{tabular}
\end{center}

\medskip
Todo lo que vive en el BMP (\lex{U+0000}--\lex{U+FFFF}) ocupa una unidad y entra sin problema; lo que
está fuera se codifica como par subrogado y no entra. La última fila es la más traicionera: es la
misma estrella de la segunda fila \emph{más} un selector de variación \lex{U+FE0F} ---lo que suelen
insertar el teclado o el celular--- y en pantalla se ve idéntica, pero son dos unidades y el literal
queda mal formado.

Conviene subrayar que esto \textbf{no} es una divergencia respecto de Java sino una coincidencia:
\lex{javac} sobre un \tok{charLiteral} con un carácter fuera del BMP responde \emph{character literal
contains more than one UTF-16 code unit}, nombrando la misma limitación por el mismo motivo. Admitir
caracteres astrales exigiría trabajar con \emph{code points} (\lex{int}) y cambiar la firma de
\lex{getNextChar()}, que es parte del material provisto.

Ésta es la primera de \textbf{dos} asimetrías entre el \tok{charLiteral} y el \tok{litString}. La
segunda, la de recuperación ante un error, está en la sección~\ref{sec:asimetria}, y las dos salen
del mismo lugar: la especificación describe cada literal con una clase de regla distinta.

\subsection{\tok{litString}}
\label{sec:string}

\begin{er}
litString = " ( Z | \bs\ C )\sups{*} "\qquad con la " de cierre NO antecedida por \bs
\medskip
Z = cualquier carácter salvo \bs, " y salto de línea
C = cualquier carácter salvo salto de línea\quad(incluye \bs\ y ")
\end{er}

\paragraph{Por qué esta ER tiene otra forma que la del \tok{charLiteral}.} La especificación
\emph{no} da una producción para el string. Da una regla \textbf{posicional}:

\begin{quote}\itshape\small
``es posible que en la secuencia de caracteres aparezca una comilla doble siempre y cuando se la
anteceda por el caracter \lex{\bs}. Por esto último la comilla doble que marca el final de un string
correcto nunca estará antecedida por ese caracter.''
\end{quote}

Es un \emph{lookbehind} de un carácter sobre el texto crudo. Consecuencia directa: \textbf{no existe
un conjunto de escapes válidos} dentro de un \tok{litString}. \lex{"hola \bs q"} es un string común y
corriente, en el que la barra y la \lex{q} son dos caracteres cualesquiera.

\begin{center}
\begin{tabular}{@{}llp{6.6cm}@{}}
\toprule
\textbf{Fuente} & \textbf{Resultado} & \textbf{Por qué} \\
\midrule
\lex{"hola \bs q"} & \tok{litString} válido & la \lex{"} de cierre no está antecedida por \lex{\bs} \\
\lex{"di \bs"hola\bs" chau"} & \textbf{un solo} \tok{litString} & las dos \lex{"} internas están antecedidas por \lex{\bs} \\
\lex{"a\bs\bs"} & \emph{sin cerrar} & la \lex{"} final está antecedida por \lex{\bs}: no cierra \\
\lex{"a\bs\bs\bs""} & \tok{litString} válido & la primera \lex{"} no cierra, la segunda sí \\
\lex{"barra \bs\bs\ al medio"} & \tok{litString} válido & la barra sólo importa pegada a la comilla de cierre \\
\bottomrule
\end{tabular}
\end{center}

\medskip
La tercera fila es el único punto en que esta regla se aparta de Java, donde \lex{"a\bs\bs"} denota
la barra invertida: aquí una barra no ``cancela'' a la anterior, porque la regla no habla de pares
sino de qué carácter antecede a la comilla de cierre. Se resolvió a favor de la especificación
escrita, que es explícita en decir \emph{nunca}.

Un \tok{litString} tampoco puede contener un salto de línea real: si aparece uno antes de la comilla
de cierre, el literal quedó sin cerrar.

\subsection{\tok{floatLiteral}}

\begin{er}
floatLiteral = \ere{Digito}{+} . \ere{Digito}{*} \ere{Exponente}{?} \ere{Sufijo}{?}
\phantom{floatLiteral }| \ere{Digito}{+} Exponente \ere{Sufijo}{?}
\phantom{floatLiteral }| \ere{Digito}{+} Sufijo
\medskip
Exponente = ( e | E ) ( + | - )\sups{?} \ere{Digito}{+}
Sufijo\phantom{nte} = f | F | d | D
\end{er}

Misma estructura que los floats de Java, con la salvedad que fija la consigna: \textbf{siempre
empiezan con un dígito}. Por eso \lex{.5} no es un float en MiniJava, sino el token \tok{punto}
seguido de \tok{intLiteral(5)}.

\paragraph{La cardinalidad del exponente, que es donde se juega todo.}\notaKleene\ En
\lex{Exponente = ( e | E ) ( + | - )}\sups{?}\lex{ }\ere{Digito}{+} conviven los tres casos, y cada
uno dice algo distinto:

\begin{description}
\item[\lex{( e | E )}] la marca de exponente es obligatoria y va \emph{sin} exponente de
  cardinalidad: aparece exactamente una vez.
\item[\lex{( + | - )}\sups{?}] el signo es \textbf{opcional} --- cero o una vez. \lex{1e5},
  \lex{1e+5} y \lex{1e-5} son los tres válidos.
\item[\ere{Digito}{+}] los dígitos del exponente son \textbf{uno o más}, \emph{nunca}
  \ere{Digito}{*}. Ésta es la diferencia que hace que \lex{1e} y \lex{1.5e+} sean errores: con
  \sups{*} el exponente vacío sería válido y \lex{1e} tokenizaría como un float.
\end{description}

En cambio, la parte fraccionaria sí lleva estrella de Kleene: \ere{Digito}{+}\lex{ . }\ere{Digito}{*}
admite \lex{1.} como float válido, con cero dígitos después del punto, igual que en Java.

Ninguna de las tres partes tiene cota superior de dígitos, a diferencia del \tok{intLiteral}. Un
float fuera del rango representable (\lex{3.5e38f}, \lex{1e999}) es un literal léxicamente
\textbf{válido}: en Java el rechazo por rango es un chequeo semántico, no léxico, y aquí se respeta
esa separación de fases.

No hay retroceso a un token más corto: una vez consumida la \lex{e} como inicio de exponente, el
literal está comprometido, y si no vienen dígitos es un error en lugar de un \tok{intLiteral}
seguido de un \tok{idMV}.

\subsection{El escape Unicode}

\begin{er}
U = \bs\ \ere{u}{+} Hex Hex Hex Hex\qquad=\qquad\bs\ \ere{u}{+} \ere{Hex}{4--4}
\end{er}

El marcador \lex{u} se repite \textbf{una o más} veces: \lex{\bs u0041}, \lex{\bs uu0041} y
\lex{\bs uuuu0041} son todos válidos y equivalentes. Es semántica heredada de Java, que la necesita
para poder representar un \lex{\bs u} literal sin re-traducirlo. Después de la última \lex{u} se
exigen exactamente cuatro dígitos hexadecimales --- no más, no menos: no es \ere{Hex}{+} ni
\ere{Hex}{*}, es una repetición acotada de extremos iguales, \ere{Hex}{4--4}.

El mismo escape sirve dentro de un \tok{charLiteral} y dentro de un \tok{litString}, y los estados
que lo reconocen están parametrizados por el estado al que deben volver.

\subsection{Puntuación}

\begin{center}
\begin{tabular}{@{}llllllll@{}}
\toprule
\textbf{Token} & \textbf{ER} & \textbf{Token} & \textbf{ER} & \textbf{Token} & \textbf{ER} & \textbf{Token} & \textbf{ER} \\
\midrule
\tok{parenA} & \lex{(} & \tok{parenC} & \lex{)} & \tok{llaveA} & \lex{\{} & \tok{llaveC} & \lex{\}} \\
\tok{corcheteA} & \lex{[} & \tok{corcheteC} & \lex{]} & \tok{puntoYComa} & \lex{;} & \tok{coma} & \lex{,} \\
\tok{punto} & \lex{.} & \tok{dosPuntos} & \lex{:} & & & & \\
\bottomrule
\end{tabular}
\end{center}

\medskip
Los diez son tokens de un solo carácter, sin lookahead: la ER de cada uno es ese carácter y nada
más.

\subsection{Operadores}

\begin{center}
\begin{tabular}{@{}llll@{}}
\toprule
\textbf{Token} & \textbf{ER} & \textbf{Token} & \textbf{ER} \\
\midrule
\tok{op\_mayor} & \lex{>} & \tok{op\_mayorIgual} & \lex{>=} \\
\tok{op\_menor} & \lex{<} & \tok{op\_menorIgual} & \lex{<=} \\
\tok{op\_negacion} & \lex{!} & \tok{op\_distinto} & \lex{!=} \\
\tok{op\_asignacion} & \lex{=} & \tok{op\_igual} & \lex{==} \\
\tok{op+} & \lex{+} & \tok{op\_incremento} & \lex{++} \\
\tok{op-} & \lex{-} & \tok{op\_decremento} & \lex{--} \\
\tok{op\_mult} & \lex{*} & \tok{op\_div} & \lex{/} \\
\tok{op\_modulo} & \lex{\%} & & \\
\tok{op\_and} & \lex{\&\&} & \tok{op\_or} & \lex{||} \\
\bottomrule
\end{tabular}
\end{center}

\medskip
Se agrupan en tres clases por su comportamiento en el autómata:

\begin{description}
\item[Con sufijo opcional] (\lex{>}, \lex{<}, \lex{!}, \lex{=}, \lex{+}, \lex{-}): el primer carácter
  ya forma un token válido y el segundo lo extiende. El estado que los reconoce es \textbf{terminal}:
  si el sufijo no viene, emite el token corto sin consumir el carácter que miró.
\item[Con segundo carácter obligatorio] (\lex{\&\&}, \lex{||}): \lex{\&} y \lex{|} sueltos \textbf{no}
  son tokens del lenguaje. El estado es \textbf{no terminal}: si el segundo carácter no viene, es un
  error léxico.
\item[De un solo carácter] (\lex{*}, \lex{\%}, y \lex{/} cuando no abre un comentario).
\end{description}

\subsection{\tok{EOF}}

Es un token más, con lexema \lex{\$}, tal como lo fija la especificación. Su ER es trivial: un solo
símbolo sobre la condición de fin de contenido.

\paragraph{El lexema \lex{\$} y el centinela de fin de archivo son cosas distintas.} El centinela que
devuelve el manejador de archivo cuando se acabó el contenido es \lex{(char) 26}, valor fijado por la
interfaz \lex{SourceManager}. El \lex{\$} es solamente el lexema que se \emph{imprime} para el token
\tok{EOF}: no se lee del fuente y no tiene ninguna relación con el centinela. Un \lex{\$} escrito en
el fuente es un símbolo inválido y produce un error léxico.

\subsubsection{Cuándo se da por terminado el archivo}

\lex{SourceManager} transporta por un mismo canal dos cosas de naturaleza distinta: los caracteres
del fuente y la señal de que los caracteres se terminaron, porque \lex{getNextChar()} devuelve un
\lex{char} y el centinela también es un \lex{char}. Eso funcionaría si el valor centinela nunca
pudiera aparecer como dato, pero \lex{0x1A} es un byte perfectamente legítimo dentro de un archivo.
Recibir el centinela es, entonces, condición \textbf{necesaria pero no suficiente}.

Por eso el fin de archivo se decide con una \textbf{conjunción}, y la interfaz se extiende con
\lex{estaEnFinDeArchivo()} para aportar el segundo término:

\begin{quote}
fin de archivo \quad$\Longleftrightarrow$\quad llegó el centinela \quad$\wedge$\quad la posición de
lectura confirma que se acabó el contenido
\end{quote}

Si llega el centinela pero la posición dice que todavía queda archivo, entonces no era un fin de
archivo: es un carácter del fuente y se reporta como símbolo inválido. Sin ese segundo término, un
\lex{0x1A} embebido cortaría el análisis en silencio, informando \lex{[SinErrores]} sobre un fuente
que nunca se terminó de leer.

En la implementación aparece un solo chequeo, porque el segundo término \emph{implica} al primero: el
manejador devuelve el centinela en una única rama, que es exactamente la misma en la que registra que
se pasó del final. La conjunción se colapsa y el $\wedge$ queda implícito; en los estados se
testea sólo el término que discrimina. La implicación no vale al revés, y ahí está todo el asunto: un
\lex{0x1A} embebido cumple el primer término y no el segundo.

\subsubsection{En qué línea cae el \tok{EOF}}

\begin{quote}
línea del \tok{EOF} \;=\; cantidad de saltos de línea consumidos \;$+$\; 1
\end{quote}

Dicho de manera más operativa: se arranca en la línea 1 y se suma 1 cada vez que se consume un
\lex{\bs n}; la línea en la que se queda parado al agotar el contenido es la del \tok{EOF}.

El caso que conviene mirar dos veces es el del archivo que \textbf{termina} en salto de línea. Un
\lex{\bs n} no es ``el final de una línea'' sino un \emph{separador}, así que después del último
salto hay una línea más ---vacía, pero existe--- y es ahí donde cae el \tok{EOF}. La consecuencia
práctica es que ese salto final tiene que entregarse al léxico como cualquier otro carácter: un
manejador que lo omita reporta el \tok{EOF} una línea antes de lo que corresponde.

La fórmula es una sola y cubre los dos casos, como muestran los tres archivos de ejemplo provistos:

\begin{center}
\begin{tabular}{@{}lccc@{}}
\toprule
\textbf{Archivo} & \textbf{Saltos} & \textbf{¿Termina en salto?} & \textbf{Línea del \tok{EOF}} \\
\midrule
\lex{lexSinErrores01.java} & 21 & sí & 22 \\
\lex{lexSinErrores02.java} & 12 & no & 13 \\
\lex{lexSinErrores03.java} & 12 & sí & 13 \\
\bottomrule
\end{tabular}
\end{center}

\subsection{Lo que no es un token}

\subsubsection{Separadores}

Espacio, tabulador y salto de línea. Se descartan sin agregarse a ningún lexema. LF, CR y CRLF
cuentan cada uno como \textbf{un} solo salto de línea, de modo que dos archivos que difieran
únicamente en el convenio de fin de línea reportan los mismos números de línea.

El \emph{form feed} y el tabulador vertical \textbf{no} son separadores en MiniJava, aunque sí lo son
en Java: no están en la lista de la especificación, y por lo tanto producen error léxico.

\subsubsection{Comentarios}

De línea (\lex{// \ldots} hasta el fin de línea) y multilínea (\lex{/* \ldots */}). Se descartan por
completo: no producen token y su contenido no forma parte de ningún lexema. \textbf{Todo} lo que está
dentro de un comentario se ignora, incluidos escapes Unicode mal formados y comillas sin cerrar.

La ER del comentario multilínea es la estándar, la misma que usan \lex{lex}/\lex{flex} para los
comentarios de C:

\begin{er}
/\bs*\quad( [\caret*] | \bs*\sups{+} [\caret*/] )\sups{*}\quad\bs*\sups{+}\quad/
\end{er}

En palabras: tras el \lex{/*}, el contenido es cualquier cantidad de (un carácter que no sea
\lex{*}) o (uno o más \lex{*} seguidos de un carácter que no sea ni \lex{*} ni \lex{/}), y cierra con
uno o más \lex{*} seguidos de \lex{/}. Los comentarios \textbf{no anidan}.

Un comentario multilínea sin cerrar es un error léxico, y se reporta con el lexema \lex{/*} y la
línea de \textbf{apertura}, no la del fin de archivo donde se descubrió.

% ===========================================================================
\section{Los autómatas}
% ===========================================================================

La implementación usa la técnica de \emph{Método por Estado}: hay una clase por estado del autómata
finito, cada una responsable de una sola transición. El driver es iterativo --- cada estado devuelve
o bien el estado siguiente, o bien el token a emitir --- de modo que la profundidad de pila no crece
con la longitud del lexema.

\subsection{Identificadores y palabras clave}

\begin{center}
\begin{tikzpicture}[shorten >=1pt,on grid,auto,
    every state/.style={draw=grisMedio,thick,fill=grisFondo,minimum size=1.05cm,font=\small},
    every edge/.style={draw=grisMedio,->,>={Stealth[round]}},
    every node/.style={font=\footnotesize}]
  \node[state,initial,initial text={}] (i)                       {$e_0$};
  \node[state,accepting]               (g) [right=3.2cm of i]    {IdGen};
  \node[state,accepting]               (c) [right=4.4cm of g]    {IdClase};
  \node[state,accepting]               (m) [below=2.6cm of g]    {IdMV};
  \path[->] (i) edge node {LetraMayus} (g)
            (g) edge node {Letra | Digito | \_} (c)
            (i) edge node[swap,pos=0.55] {LetraMin} (m);
  \path[->] (c) edge [loop above] node {Letra | Digito | \_} ();
  \path[->] (m) edge [loop below] node {Letra | Digito | \_} ();
\end{tikzpicture}
\end{center}

\noindent
El estado \textsf{IdGen} es \textbf{terminal}: si el token termina ahí, es \tok{idGen}; si aparece un
carácter más, transiciona a \textsf{IdClase}. Al emitir desde \textsf{IdMV} se consulta la tabla de
lexemas: si el lexema completo coincide con una palabra reservada, se emite la palabra clave.

\subsection{\tok{charLiteral}}

\begin{center}
\begin{tikzpicture}[shorten >=1pt,node distance=3.1cm,on grid,auto,
    every state/.style={draw=grisMedio,thick,fill=grisFondo,minimum size=1.0cm,font=\small},
    every edge/.style={draw=grisMedio,->,>={Stealth[round]}},
    every node/.style={font=\small}]
  \node[state,initial,initial text={}] (i)                {$e_0$};
  \node[state]                         (a) [right=of i]   {Abierto};
  \node[state]                         (e) [below=2.3cm of a] {Escapado};
  \node[state]                         (c) [right=3.6cm of a] {Cerrando};
  \node[state,accepting]               (f) [right=3.4cm of c] {char};
  \path[->] (i) edge node {\lex{'}} (a)
            (a) edge node {$\neq$ \lex{\bs} , \lex{'}} (c)
            (a) edge node[swap] {\lex{\bs}} (e)
            (e) edge node[swap,pos=0.4] {$\neq$ \lex{u}} (c)
            (c) edge node {\lex{'}} (f);
\end{tikzpicture}
\end{center}

\noindent
Desde \textsf{Escapado}, una \lex{u} deriva al autómata del escape Unicode, que vuelve a
\textsf{Cerrando} tras consumir los cuatro dígitos hexadecimales. Los tres caminos que no llegan al
estado final son los tres errores del \tok{charLiteral}: \lex{''} (vacío), un fin de línea o de
archivo antes de cerrar (\emph{sin cerrar}), y cualquier otro carácter donde debía ir la comilla de
cierre (\emph{mal formado}).

\subsection{\tok{litString}}

El \emph{lookbehind} de la sección~\ref{sec:string} es un bit de memoria, y se modela con dos
estados: en \textsf{Abierto} el carácter anterior no fue una barra; en \textsf{TrasBarra} sí lo fue.

\begin{center}
\begin{tikzpicture}[shorten >=1pt,on grid,auto,
    every state/.style={draw=grisMedio,thick,fill=grisFondo,minimum size=1.0cm,font=\small},
    every edge/.style={draw=grisMedio,->,>={Stealth[round]}},
    every node/.style={font=\footnotesize}]
  \node[state,initial,initial text={}] (i)                    {$e_0$};
  \node[state]                         (a) [right=2.7cm of i] {Abierto};
  \node[state]                         (b) [right=6.2cm of a] {TrasBarra};
  \node[state,accepting]               (f) [below=3.0cm of a] {litString};
  \path[->] (i) edge node {\lex{"}} (a);
  \path[->] (a) edge [bend left=22] node[above,pos=0.5] {\lex{\bs}} (b);
  \path[->] (b) edge [bend left=22] node[below,pos=0.5] {cualquier otro, incluida \lex{"}} (a);
  \path[->] (a) edge node[swap] {\lex{"}} (f);
  \path[->] (a) edge [loop above] node[above] {$\neq$ \lex{\bs} , \lex{"} , salto} ();
  \path[->] (b) edge [loop above] node[above] {\lex{\bs}} ();
\end{tikzpicture}
\end{center}

\noindent
Hay que reparar en el bucle de \textsf{TrasBarra} sobre \lex{\bs}: una barra que sigue a otra barra
\textbf{no} devuelve al estado \textsf{Abierto}, porque la comilla que venga después seguirá estando
antecedida por una barra. Eso es exactamente lo que hace que \lex{"a\bs\bs"} quede sin cerrar. Desde
\textsf{TrasBarra}, una \lex{u} deriva al escape Unicode.

\subsection{\tok{floatLiteral}}

\begin{center}
\begin{tikzpicture}[shorten >=1pt,on grid,auto,
    every state/.style={draw=grisMedio,thick,fill=grisFondo,minimum size=0.95cm,font=\small},
    every edge/.style={draw=grisMedio,->,>={Stealth[round]}},
    every node/.style={font=\footnotesize}]
  \node[state,initial,initial text={}] (i)                     {$e_0$};
  \node[state,accepting]               (e) [right=2.9cm of i]  {Entero};
  \node[state,accepting]               (p) [right=4.2cm of e]  {ConPunto};
  \node[state]                         (x) [below=3.2cm of i]  {ExpInicio};
  \node[state]                         (r) [right=3.9cm of x]  {ReqDigito};
  \node[state,accepting]               (d) [right=4.2cm of r]  {ExpDigitos};
  \path[->] (i) edge node {Digito} (e);
  \path[->] (e) edge node {\lex{.}} (p);
  % Las dos entradas al exponente llegan por lados distintos para que sus
  % etiquetas no compartan espacio.
  \path[->] (e) edge node[swap,pos=0.55] {\lex{e} | \lex{E}} (x);
  \path[->] (p) edge node[pos=0.22,right=2pt] {\lex{e} | \lex{E}} (x);
  \path[->] (x) edge node[above] {\lex{+} | \lex{-}} (r);
  \path[->] (x) edge [bend right=22] node[below,pos=0.55] {Digito} (d);
  \path[->] (r) edge node[above] {Digito} (d);
  \path[->] (e) edge [loop above] node[above] {Digito} ();
  \path[->] (p) edge [loop above] node[above] {Digito} ();
  \path[->] (d) edge [loop below] node[below] {Digito} ();
\end{tikzpicture}
\end{center}

\noindent
\textsf{Entero}, \textsf{ConPunto} y \textsf{ExpDigitos} son terminales; \textsf{ExpInicio} y
\textsf{ReqDigito} no lo son, y ahí está representado el \ere{Digito}{+} del exponente: llegar a
\textsf{ExpInicio} y no encontrar un dígito ni un signo es el error de \lex{1e}. Los sufijos
\lex{f}/\lex{F}/\lex{d}/\lex{D} pueden cerrar desde \textsf{Entero}, \textsf{ConPunto} o
\textsf{ExpDigitos}; se omiten del diagrama para no saturarlo.

Los bucles sobre \textsf{Entero}, \textsf{ConPunto} y \textsf{ExpDigitos} son correctos acá y no
contradicen la cadena de nueve estados del \tok{intLiteral}: ninguna de las tres partes de un
\tok{floatLiteral} tiene cota de dígitos, así que la repetición \emph{sí} es una estrella de Kleene.
La cota de nueve pertenece únicamente al literal entero.

\subsection{Comentario multilínea}

\begin{center}
\begin{tikzpicture}[shorten >=1pt,node distance=4.0cm,on grid,auto,
    every state/.style={draw=grisMedio,thick,fill=grisFondo,minimum size=1.0cm,font=\small},
    every edge/.style={draw=grisMedio,->,>={Stealth[round]}},
    every node/.style={font=\small}]
  \node[state,initial,initial text={\lex{/*}}] (s)              {SinAst};
  \node[state]                                 (c) [right=of s] {ConAst};
  \node[state,accepting]                       (f) [right=of c] {cierra};
  \path[->] (s) edge [bend left=12] node {\lex{*}} (c)
            (c) edge [bend left=12] node {$\neq$ \lex{*} , \lex{/}} (s)
            (c) edge node {\lex{/}} (f);
  \path[->] (s) edge [loop above] node {$\neq$ \lex{*}} ();
  \path[->] (c) edge [loop above] node {\lex{*}} ();
\end{tikzpicture}
\end{center}

\subsection{El reparto entre autómata y tabla}
\label{sec:tabla}

De los 56 tokens, \textbf{49 tienen lexema fijo}, en el sentido de que el tipo y el lexema se
determinan mutuamente: las 21 palabras clave, los 10 símbolos de puntuación, los 17 operadores y
\tok{EOF}. Los otros 7 --- los tres identificadores y los cuatro literales --- leen su lexema del
fuente.

Los 49 no necesitan un estado que lleve su tipo encima: alcanza con que el autómata delimite el
lexema y con consultar una tabla. Lo único que la tabla \emph{no} puede resolver es el
\emph{maximal munch}: no sabe si un \lex{>} es \tok{op\_mayor} o el principio de un
\tok{op\_mayorIgual}, y eso lo decide el autómata mirando el carácter siguiente. El reparto queda
nítido:

\begin{quote}
\textbf{el autómata delimita el lexema; la tabla lo nombra.}
\end{quote}

Es la técnica recomendada para las palabras clave, generalizada a los 49. Modelar cada palabra
reservada con estados dedicados exigiría casi un estado por letra, fusionados con los estados de
identificador: decenas de estados frágiles para algo que una tabla resuelve en tiempo constante.

% ===========================================================================
\section{Decisiones tomadas}
% ===========================================================================

Dos reglas gobiernan todas las decisiones del proyecto:

\begin{enumerate}[leftmargin=1.6em]
\item \textbf{Lo que la especificación fija, se respeta al pie de la letra.}
\item \textbf{Lo que no fija, se resuelve como en Java} --- regla escrita explícitamente en la
  consigna.
\end{enumerate}

\medskip
Los casos que admitían más de una lectura, y cómo se resolvió cada uno:

\begin{center}
\begin{longtable}{@{}p{4.3cm}p{5.0cm}p{5.6cm}@{}}
\toprule
\textbf{Caso} & \textbf{Decisión} & \textbf{Fundamento} \\
\midrule
\endhead
\lex{float} y \lex{double} como palabras reservadas &
\lex{float} es reservada; \lex{double} \textbf{no}: sigue siendo un \tok{idMV} común &
\lex{float} se agregó por confirmación directa de la cátedra; que MiniJava tenga un único tipo de
punto flotante no convierte a \lex{double} en reservada \\

Gramática del \tok{floatLiteral} &
la de Java, con la salvedad de que siempre empieza con dígito &
la consigna la describe en una sola frase: ``la misma estructura que los floats de Java'' \\

Límite de dígitos de los floats &
sin límite, a diferencia del \tok{intLiteral} &
la cota de 9 dígitos está enunciada sólo para el literal entero \\

Float fuera de rango (\lex{1e999}) &
literal léxicamente válido &
en Java el rechazo por rango es un chequeo \emph{semántico}; el léxico no valida magnitudes \\

Escapes en \tok{charLiteral} &
cualquier carácter tras la barra &
la especificación da una producción con dos alternativas y ``cualquier caracter'' explícito \\

Escapes en \tok{litString} &
no hay conjunto de escapes: cierra la primera \lex{"} no antecedida por \lex{\bs} &
la especificación da una regla \emph{posicional}, no una producción; ver sección~\ref{sec:string} \\

\lex{\&} y \lex{|} sueltos &
error léxico &
la lista de operadores incluye \lex{\&\&} y \lex{||}, y no los incluye sueltos \\

\lex{''} (char vacío) &
error léxico &
no está en ninguna de las dos alternativas de la producción \\

Marcador \lex{u} repetido &
\lex{\bs uu0041} es válido y equivale a \lex{\bs u0041} &
semántica de Java (JLS §3.3), verificada contra la especificación \\

Maximal munch &
se consume siempre el token más largo posible, sin retroceso &
comportamiento estándar; \lex{1234abc} da \tok{intLiteral(1234)} y \tok{idMV(abc)} \\

Escape Unicode mal formado en un \tok{litString} &
se reporta y se descarta el literal hasta su comilla de cierre, sin emitir token &
sobre \lex{"a\bs u12G4" + v1} se emite un error y luego \tok{op+} e \tok{idMV(v1)}: el literal roto
no se emite, y el análisis retoma en el token siguiente \\

Rango de un \tok{intLiteral} &
la cota es de \emph{dígitos}, no de valor: \lex{0000000001} es error &
la especificación acota dígitos, no magnitud \\
\bottomrule
\end{longtable}
\end{center}

% ===========================================================================
\section{Manejo de errores}
% ===========================================================================

\subsection{Formato del reporte}

Ante un error, el léxico lanza una excepción que el Módulo Principal captura y formatea. La última
línea es siempre el código de contrato \lex{[Error:Lexema|NroLinea]}; las anteriores son el reporte
legible, que es uno de los logros perseguidos.

\begin{codigoLlano}
(idMV,v1,2)
Error Léxico en línea 2, columna 6: # no es un símbolo válido
Detalle: v1 + # chau
              ^

[Error:#|2]
(idMV,chau,2)
\end{codigoLlano}

\textbf{El lexema reportado es siempre el armado hasta el momento del error}, como exige la
consigna, y el código de contrato lleva ese mismo lexema. Si el carácter ofensivo \emph{es} el
problema, queda incluido en el lexema; si el problema es el lexema ya construido (un \lex{\&} sin su
par, un entero de más de nueve dígitos, un exponente sin dígitos), el carácter que se está mirando es
un token legítimo que no tiene nada que ver con el error, y por lo tanto no se consume ni se incluye.

Los lexemas de error nunca llevan un carácter invisible crudo: un salto de línea dentro de un lexema
partiría la línea del reporte en dos y rompería el formato de un elemento por línea. Se renderizan
con su escape (\lex{\bs n}, \lex{\bs t}, \lex{\bs u00a0}).

\subsection{Criterio de recuperación}

\begin{quote}
\textbf{Todo estado que lanza un error consume del fuente exactamente los caracteres que reporta en
el lexema, ni uno más.}
\end{quote}

No hay un paso de recuperación posterior: el siguiente pedido de token arranca en el primer carácter
que \textbf{no} formó parte del lexema reportado. Esto coincide con lo que hace \lex{javac}, que
tampoco descarta caracteres ajenos al lexema ofensivo, y garantiza estructuralmente la terminación:
el lexema de un error nunca es vacío, así que cada error consumió al menos un carácter y la posición
de lectura siempre avanza.

La única excepción es el escape Unicode mal formado dentro de un \tok{litString}, que además descarta
lo que queda del literal hasta su comilla de cierre. Es recuperación en modo pánico con token de
sincronización: el delimitador es la comilla de cierre, y el estado conocido al que se vuelve es
``fuera de todo literal''. Sin resincronizar, la comilla de \emph{cierre} abriría un literal nuevo
que se comería el resto de la línea.

\subsection{Por qué el \tok{charLiteral} no resincroniza}
\label{sec:asimetria}

Que el \tok{litString} descarte hasta su comilla de cierre y el \tok{charLiteral} no descarte nada no
es una inconsistencia: es que el error deja al analizador en lugares distintos.

\begin{quote}
Cuando un \tok{litString} falla, el análisis sigue \textbf{adentro} del literal, y lo que viene es
contenido de largo desconocido ---texto, no código---. Descartarlo descarta contenido.

Cuando un \tok{charLiteral} falla, el análisis está en un punto donde el literal \textbf{ya debería
haber terminado}: su contenido es exactamente un carácter. Lo que sigue no hay razón para creer que
sea contenido del literal. Descartarlo descartaría código.
\end{quote}

De ahí salen tres consecuencias concretas:

\begin{description}
\item[No hay nada que descartar.] El modo pánico existe para saltear una corrida de basura de largo
  arbitrario. En un \tok{charLiteral} esa corrida tiene largo a lo sumo uno: al detectar el error se
  consumieron la comilla de apertura y un carácter.
\item[La comilla simple no es un delimitador confiable.] En un \tok{litString} la \lex{"} de cierre
  es inequívoca, porque todo lo que está entre comillas es contenido. En un \tok{charLiteral}, una
  \lex{'} posterior puede ser tanto el cierre de éste como la apertura del siguiente ---y muchas
  veces no hay ninguna: en \lex{'a; x = 1;} no aparece una segunda comilla en toda la línea, así que
  descartar hasta el cierre se comería cinco tokens legítimos.
\item[La recuperación estructural ya alcanza.] Como el literal es acotado, el criterio general
  ---consumir exactamente lo que se reporta--- basta por sí solo: al reportar \lex{'a} la lectura
  quedó sobre el carácter ofensivo y el análisis retoma ahí sin ningún paso adicional.
\end{description}

\paragraph{El costo, asumido a propósito.} Un \lex{'abc'} produce \emph{dos} errores y un
\tok{idMV(bc)} de más, en vez de un único error limpio:

\begin{codigoLlano}
'ab'; var x = 1;
  -> Error Léxico en línea 1, columna 1: charLiteral mal formado, contiene más de un carácter: 'a
     [Error:'a|1]
     (idMV,b,1)
     Error Léxico en línea 1, columna 4: charLiteral mal formado, contiene más de un carácter: ';
     [Error:';|1]
     (pr_var,var,1) (idMV,x,1) (op_asignacion,=,1) (intLiteral,1,1) (puntoYComa,;,1)
\end{codigoLlano}

\paragraph{Es idéntico a lo que hace \lex{javac}.} No es una analogía: sobre el mismo caso, la
recuperación de \lex{javac} coincide token por token con la de este compilador. Su salida real,
verificada ejecutándolo:

\begin{codigoLlano}
$ javac T.java
T.java:1: error: unclosed character literal
class T { void m() { char c = 'ab'; int x = 1; } }
                              ^
T.java:1: error: unclosed character literal
class T { void m() { char c = 'ab'; int x = 1; } }
                                 ^
T.java:1: error: not a statement
class T { void m() { char c = 'ab'; int x = 1; } }
                                ^
3 errors
\end{codigoLlano}

\noindent
La correspondencia es exacta en las cuatro decisiones que definen la recuperación:

\begin{center}
\begin{tabular}{@{}p{6.1cm}ll@{}}
\toprule
\textbf{Decisión} & \lex{javac} & \textbf{este compilador} \\
\midrule
Errores léxicos sobre \lex{'ab'} & 2 & 2 \\
Dónde caen & en las dos comillas & en las dos comillas \\
Qué pasa con la \lex{b} del medio & identificador & \tok{idMV(b)} \\
Qué pasa con el código posterior & sobrevive entero & sobrevive entero \\
\bottomrule
\end{tabular}
\end{center}

\medskip
En \lex{javac} la tercera línea ---\emph{not a statement}, con el \lex{\caret} sobre la \lex{b}--- es
ya un error del análisis sintáctico, y es justamente la prueba de que su léxico entregó la \lex{b}
como un identificador en vez de descartarla. La única diferencia está en el \emph{texto} del
diagnóstico, no en la recuperación: \lex{javac} usa un mismo mensaje para los dos errores, mientras
que acá se distingue \emph{sin cerrar} ---lo cortó un fin de línea o de archivo--- de \emph{mal
formado} ---contiene más de un carácter---, que son dos problemas distintos y ayudan distinto.

\paragraph{Por qué se acepta ese costo.} Porque la alternativa es peor. Lo que se ve en la última
línea del ejemplo es que todo el código posterior al literal roto sobrevive intacto; si se descartara
hasta la comilla de cierre, se perdería. Se prefiere un error de más antes que un token legítimo de
menos: un error extra el lector lo ignora, mientras que un token que desaparece arruina el
diagnóstico de todo lo que viene después.

\subsection{Dónde apunta el \lex{\caret}}

Se adoptó el criterio de \lex{javac}, verificado ejecutándolo sobre cada caso:

\begin{quote}
Si el mensaje habla de un \textbf{carácter ofensivo puntual}, el \lex{\caret} va a ese carácter. Si habla
de la \textbf{construcción entera} que falló, el \lex{\caret} va a donde esa construcción \textbf{empieza}.
\end{quote}

\subsection{Ubicar el error sin releer el archivo}

El reporte necesita la columna y el texto completo de la línea del error. Ninguna de las dos cosas
obliga a una segunda pasada sobre el fuente:

\begin{center}
\begin{tabular}{@{}lll@{}}
\toprule
\textbf{Operación} & \textbf{Costo} & \textbf{Cómo} \\
\midrule
columna & $O(1)$ & contador que avanza al leer y se reinicia en cada salto de línea \\
texto de la línea & $O(\text{largo de la línea})$ & otro contador guarda dónde empieza la línea actual \\
\bottomrule
\end{tabular}
\end{center}

\medskip
La información de posición se construye \emph{mientras} se lee, con dos enteros, y el recorrido para
recuperar la línea está acotado por su largo, que es exactamente el texto que el reporte va a
imprimir de todos modos.

% ===========================================================================
\section{Logros perseguidos}
% ===========================================================================

El proyecto aspira a \textbf{los ocho logros} de la Etapa 1, perseguidos desde el comienzo del
desarrollo y no como agregados posteriores sobre un mínimo producto viable:

\begin{enumerate}[leftmargin=1.8em,itemsep=0.3ex]
\item Floats I
\item Unicodes
\item Multi-detección de Errores Léxicos
\item Reporte de Error Elegante
\item Manejador de Archivos Eficiente
\item Imbatibilidad Léxica
\item Good Job! I
\item Entrega Anticipada Léxica
\end{enumerate}

Cuatro de ellos ---Floats I, Unicodes, Multi-detección y Reporte de Error Elegante--- cambian el
comportamiento observable del compilador. Para que tanto el comportamiento base que pide la consigna
como el del logro sean demostrables, cada uno vive detrás de un interruptor de configuración, todos
en \lex{true} por omisión:

\begin{codigoLlano}
ConfiguracionDelCompilador(floatsHabilitados, unicodeHabilitado,
                           reporteElegante, multiDeteccionDeErrores)
\end{codigoLlano}

Sin esos interruptores el comportamiento base quedaría sepultado bajo el del logro y no habría forma
de mostrar los dos.

% ===========================================================================
\section{Estrategia de prueba}
% ===========================================================================

El catálogo de casos de test \textbf{es el propio directorio de casos} del repositorio, con el
oráculo escrito en las primeras líneas de cada archivo en el formato que definen los testers
provistos. Como esas líneas son comentarios MiniJava válidos, el archivo es a la vez el caso de
prueba y su propia especificación: el compilador las ignora y el tester las lee.

\begin{codigoLlano}
resources/
    sinErrores/    76 archivos   oráculo: //#lexema,nroLinea
    conErrores/    82 archivos   oráculo: // [Error:lexema|nroLinea]
    regresion/                   entradas que rompieron una corrida de fuzzing
\end{codigoLlano}

\subsection{Cómo se declara el oráculo de un caso}

Un caso de \lex{sinErrores/} declara, con una línea \lex{//\#lexema,nroLinea} por token, la secuencia
completa que debe salir. Un caso de \lex{conErrores/} declara además, con
\lex{// [Error:lexema|nroLinea]}, cada código de error esperado.

El tester provisto verifica únicamente que el \emph{primer} código de error declarado aparezca en
algún lugar de la salida. Eso deja dos huecos: un caso con varios errores queda casi sin verificar, y
nada comprueba qué tokens salen entre medio ---que es justamente donde se esconderían los errores de
recuperación---. Por eso la cabecera admite una marca extra:

\begin{quote}
\lex{// [Completo]} --- de esta línea en adelante, la secuencia declarada debe coincidir
\textbf{exactamente} con la salida real: todos los códigos de error y todos los tokens, en orden y
sin nada de más ni de menos.
\end{quote}

Sin esa marca la verificación es de \emph{subsecuencia}: los elementos declarados tienen que
aparecer todos y en orden, pero se admite que haya otros entre medio. Con la marca es de igualdad.
Un caso completo se ve así:

\begin{codigoLlano}
// [Error:"a\u12|6]     <- primera línea: el código que lee el tester provisto
// [Completo]           <- de acá en más, verificación exacta
// [Error:"a\u12|6]
//#+,6
//#v1,6
//#$,7
"a\u12G4" + v1
\end{codigoLlano}

\noindent
Se lee así: sobre ese fuente, el compilador debe emitir el error del escape Unicode mal formado con
lexema \lex{"a\bs u12} en la línea 6, después el token \lex{+}, después \lex{v1}, y cerrar con el
\tok{EOF} en la línea 7 --- exactamente eso y nada más. La primera línea está duplicada a propósito:
el tester provisto le hace \lex{substring(3)} a la línea 1 del archivo, así que ahí tiene que estar
sí o sí el código de error.

% ===========================================================================
\section{Divergencias deliberadas respecto de Java}
\label{sec:divergencias}
% ===========================================================================

Ninguna de éstas es un descuido: todas están cubiertas por casos de test escritos a propósito, y
todas tienen su razón.

\begin{center}
\begin{longtable}{@{}p{3.4cm}p{4.6cm}p{6.4cm}@{}}
\toprule
\textbf{Caso} & \textbf{MiniJava} & \textbf{Por qué} \\
\midrule
\endhead
\lex{\bs uXXXX} fuera de un literal &
en un comentario es texto común &
en Java el escape se traduce antes de tokenizar, en todo el archivo; implementar ese
pre-procesamiento es una fase entera aparte que la consigna no pide \\

\lex{041}, \lex{0x1F}, \lex{0b1010} &
decimal; no hay literales octales, hexadecimales ni binarios &
la ER de la especificación es decimal y nada más \\

\lex{1\_000} &
\tok{intLiteral(1)}, error por el \lex{\_}, \tok{intLiteral(000)} &
ídem: no hay separadores de miles \\

\lex{0000000001} &
error &
la cota es de \emph{dígitos}, no de valor como en Java \\

\lex{1e999} &
\tok{floatLiteral} válido &
el rango es un chequeo semántico, no léxico \\

\lex{.5} &
\tok{punto} + \tok{intLiteral} &
la consigna exige que un float empiece con dígito \\

\lex{'\bs q'} &
\tok{charLiteral} válido &
permisividad explícita de la especificación \\

\lex{"a\bs q"} &
\tok{litString} válido &
la especificación no define un conjunto de escapes: define qué comilla cierra \\

\lex{"a\bs\bs"} &
\emph{sin cerrar} &
la regla es \textbf{posicional}: la comilla de cierre nunca está antecedida por \lex{\bs}, y una
barra no cancela a la anterior \\

Dígitos árabe-índicos &
son un \tok{intLiteral} &
\lex{Character.isDigit} los acepta; restringir a ASCII contradiría el logro Unicodes \\

Un ideograma CJK (\lex{U+4E2D}) al empezar un identificador &
error, aunque sí puede continuarlo &
la ER dice ``letra mayúscula'' / ``letra minúscula'', y no es ninguna de las dos \\

\emph{Form feed} y tabulador vertical &
error léxico &
no están en la lista de separadores \\

BOM al inicio del archivo &
carácter desconocido &
ídem \\
\bottomrule
\end{longtable}
\end{center}

\end{document}
